\chapter{Graph}

\section{Fundamentals}
	\kactlimport{BellmanFord.h}
	\kactlimport{FloydWarshall.h}
	\kactlimport{TopoSort.h}
	\kactlimport{MinMeanCycle.h}
	\kactlimport{SteinerTree.h}

\section{Network flow}
	\kactlimport{PushRelabel.h}
	\kactlimport{Dinic.h}
	\kactlimport{MinCostMaxFlow.h}
	\kactlimport{MinCut.h}
	\kactlimport{GlobalMinCut.h}
	\kactlimport{GomoryHu.h}

\section{Matching}
	\kactlimport{HopcroftKarp.h}
	\kactlimport{DFSMatching.h}
	\kactlimport{MinimumVertexCover.h}
	\kactlimport{WeightedMatching.h}
	\kactlimport{GeneralMatching.h}

\section{DFS algorithms}
	\kactlimport{SCC.h}
	\kactlimport{BiconnectedComponents.h}
	\kactlimport{2sat.h}
	\kactlimport{EulerWalk.h}
	\kactlimport{DominatorTree.h}
	\kactlimport{DynamicConnectivity.h}

\section{Coloring}
	\kactlimport{EdgeColoring.h}

\section{Trees}
	\kactlimport{BinaryLifting.h}
	\kactlimport{LCA.h}
	\kactlimport{CompressTree.h}
	\kactlimport{HLD.h}
	\kactlimport{CentroidDecomposition.h}
	\kactlimport{LinkCutTree.h}
	\kactlimport{DirectedMST.h}

\section{Math}
	\subsection{Number of Spanning Trees}
		% I.e. matrix-tree theorem.
		% Source: https://en.wikipedia.org/wiki/Kirchhoff%27s_theorem
		% Test: stress-tests/graph/matrix-tree.cpp
		Create an $N\times N$ matrix \texttt{mat}, and for each edge $a \rightarrow b \in G$, do
		\texttt{mat[a][b]--, mat[b][b]++} (and \texttt{mat[b][a]--, mat[a][a]++} if $G$ is undirected).
		Remove the $i$th row and column and take the determinant; this yields the number of directed spanning trees rooted at $i$
		(if $G$ is undirected, remove any row/column). Determinant mod $p$: IntDeterminant.

	\subsection{Erdős–Gallai theorem}
		% Source: https://en.wikipedia.org/wiki/Erd%C5%91s%E2%80%93Gallai_theorem
		% Test: stress-tests/graph/erdos-gallai.cpp
		A simple graph with node degrees $d_1 \ge \dots \ge d_n$ exists iff $d_1 + \dots + d_n$ is even and for every $k = 1\dots n$,
		\[ \sum _{i=1}^{k}d_{i}\leq k(k-1)+\sum _{i=k+1}^{n}\min(d_{i},k). \]

	\subsection{Matchings, covers, flows}
		\label{sec:konig}
		\textbf{K\H{o}nig} (bipartite): max matching = min vertex cover;
		max independent set = $V$ $-$ max matching; min edge cover = $V$ $-$ max matching
		(no isolated vertices). Min vertex cover from a max flow: left vertices \emph{not} reachable
		from $s$ in the residual graph + right vertices reachable (MinimumVertexCover).\\
		\textbf{Hall:} a left-perfect matching exists iff $|N(S)| \ge |S|$ for all left $S$.\\
		\textbf{DAG path cover:} min vertex-disjoint path cover = $V$ $-$ max matching on (out-copy, in-copy);
		with paths allowed to share vertices: first take the transitive closure.
		\textbf{Dilworth:} min chains = max antichain.\\
		\textbf{Max-flow = min-cut.} Lower bounds $l \le f \le c$ on $u \to v$: give it capacity $c - l$ and
		add $S' \to v$ and $u \to T'$ with capacity $l$ (new source/sink), plus $t \to s$ with $\infty$;
		feasible iff the max flow $S' \to T'$ saturates all $S'$ edges. Then max flow $s \to t$ in what is left.\\
		\textbf{Project selection / closure:} profit $p_i > 0$: $s \to i$ ($p_i$); cost: $i \to t$ ($-p_i$);
		``$i$ requires $j$'': $i \to j$ ($\infty$). Best $= \sum_{p_i > 0} p_i -$ mincut.
